\documentclass[10pt,a4paper]{amsart}
\usepackage{amsmath,amssymb,amsthm,graphicx}
\usepackage{amsmath}
\usepackage{amssymb}
\usepackage{hyperref}
\newtheorem{Theorem}{Theorem}[section]
\newtheorem{Corollary}[Theorem]{Corollary}
\newtheorem{Lemma}[Theorem]{Lemma}
\newtheorem{Proposition}[Theorem]{Proposition}
{\theoremstyle{definition}
\newtheorem{Note}[Theorem]{Note}
\newtheorem{Definition}[Theorem]{Definition}
\newtheorem{Problem}[Theorem]{Problem}
}
\usepackage{cleveref}
\crefname{originaldefinition}{original definition}{original definitions}
\Crefname{originaldefinition}{Original definition}{Original definitions}
\makeatletter
\newlabel{def:ps}{{1.14}{source}{}{sourcepointer-1}{}}
\newlabel{def:ps@cref}{{[originaldefinition][0][]1.14}{[1][0][]source}{}{}{}}
\makeatother
\crefname{originallemma}{original lemma}{original lemmas}
\Crefname{originallemma}{Original lemma}{Original lemmas}
\makeatletter
\newlabel{lem:4}{{1.15}{source}{}{sourcepointer-2}{}}
\newlabel{lem:4@cref}{{[originallemma][0][]1.15}{[1][0][]source}{}{}{}}
\makeatother
\title{The action of PSL\_2(Z) on the algebra Lambda is faithful}
\author{Paul Terwilliger}
\date{doi:10.3842/SIGMA.2011.069; CC BY4.0}
\begin{document}
\maketitle
\noindent\textit{Editorial note.} All delivered mathematical text is the author's own, copied from the pinned publisher source sources/sigma-2011-069-askey-wilson/pinned-original-source.tex; no statement or proof step was rewritten, reordered or removed.\par\medskip
\noindent\textit{Editorial note.} Preamble provenance: the retained preamble is the source's own layout adapter plus the author's macro block at lines 10-18 of sources/sigma-2011-069-askey-wilson/pinned-original-source.tex (everything between the publisher class metadata and \textbackslash{}begin\{document\}, comments excluded). Nothing was added, no macro was invented, and the publisher document class is deliberately not reproduced; the excerpt serves only as preamble, never as a span coordinate.\par\medskip
\noindent\textit{Editorial note.} External references are navigation pointers only; every deliberately kept step of the author's argument stays inside the counted spans. The displayed numbers of those pointers are the author's own numbering under the author's \textbackslash{}newtheorem declarations (the source has no subsections, so numbering is section.index), computed mechanically by the lane's scanner.\par\medskip
\noindent\textit{Editorial note.} Two references to other labelled items occur inside the span, Definition def:ps and Lemma lem:4(i); both are recorded as navigation stubs with the authors' own numbers and the correct kind, and every computation that consumes them stays inside the counted span.\par\medskip
\noindent\textit{Editorial note.} The last line of the proof ('the result follows.') is the author's own closing, kept verbatim.\par\medskip
\noindent\textit{Editorial note.} Licence binding: the card's source\_version now carries the value recorded in this source's own licence evidence (sources/sigma-2011-069-askey-wilson/license-observation.json), namely 'SIGMA7(2011),069; publisher final form2011-07-09'. The adopt step of the publication line compares the two and refused the earlier revision; the card is aligned to the evidence rather than the evidence to the card.\par\medskip
\noindent\textit{Allowed standard prerequisites (editorial).}\begin{itemize}
\item The matrices p and s generating the action, with s of the form s\_0 + s\_1 lambda. Same source, Definition labelled def:ps; invoked inside the counted proof at lines 858-883. Lets a word in s and p\textasciicircum{}\{+-1\} be read as a polynomial in lambda with coefficients in 2x2 matrices.
\item The determinant and trace formulas for the elements of Lambda. Same source, Lemma labelled lem:4, clause (i); invoked inside the counted proof at lines 900-903. Converts the centrality equation eta = kappa I into the identity (-lambda)\textasciicircum{}ell = kappa\textasciicircum{}2, which forces ell to be even.
\end{itemize}\medskip
\ifdefined\hypertarget\hypertarget{sourcepointer-1}{}\fi
\noindent\textit{Source reference.} Original definition 1.14: sources/sigma-2011-069-askey-wilson/pinned-original-source.tex, source lines 858--858. Navigation pointer for \textbackslash{}ref\{def:ps\} used inside the counted span (line 858); the argument that consumes it is retained in the span.\par\medskip
\ifdefined\hypertarget\hypertarget{sourcepointer-2}{}\fi
\noindent\textit{Source reference.} Original lemma 1.15: sources/sigma-2011-069-askey-wilson/pinned-original-source.tex, source lines 902--902. Navigation pointer for \textbackslash{}ref\{lem:4\} used inside the counted span (line 902); the argument that consumes it is retained in the span.\par\medskip
\begin{Lemma}
\label{lem:faith1}
The action of
${\rm  {PSL}}_2(\mathbb Z)$ on $\Lambda $ is faithful.
\end{Lemma}

\begin{proof}
Pick an integer $n\geq 1$.
Consider an element $\eta \in \Lambda $
of the form
$\eta =\eta_1 \eta_2 \cdots \eta_n$ such that
 for $1 \leq i \leq n$,  $\eta_i=s$
for one parity of $i$ and
$\eta_i \in \lbrace p,p^{-1} \rbrace$ for the other parity of~$i$.
We show that $\eta \not\in Z(\Lambda)$.
To this end we assume
 $\eta \in Z(\Lambda )$ and get a contradiction.
Let $\ell$ denote the number of times~$s$ occurs among
$\lbrace \eta_i\rbrace_{i=1}^n$.
Assume for the moment $\ell = 0$.
Then $n=1$ so
$\eta =\eta_1 \in \lbrace p, p^{-1}\rbrace$.
The elements $p$, $p^{-1}$ are not in~$Z(\Lambda)$, for a contradiction.
Therefore $\ell \neq 0$.
From the nature of the matrices~$p$,~$s$ in Def\/inition~\ref{def:ps},
we may
view $\eta $ as a polynomial in $\lambda$
that has
coef\/f\/icients in
${\rm {Mat}}_2(\mathbb F)$ and degree at most $\ell$.
Call this polynomial~$f$.
We claim that the degree of~$f$ is exactly~$\ell$.
To prove the claim,
write
\begin{gather*}
s = s_0+ s_1 \lambda,   \qquad
s_0= \left(
\begin{array}{  c c}
0 & 1  \\
0 & 0
\end{array}
\right),
\qquad
s_1= \left(
\begin{array}{  c c}
0 & 0  \\
1 & 0
\end{array}
\right).
\end{gather*}
Let $m
\in {\rm {Mat}}_2(\mathbb F)$ denote the coef\/f\/icient of
$\lambda^\ell$ in $f$.
The matrix $m$ is obtained from
$\eta_1 \eta_2\cdots \eta_n$ by
replacing each occurrence of $s$ by $s_1$.
Using $s_1p s_1=-s_1$ and
 $s_1p^{-1} s_1=s_1$  we f\/ind
$m\in \lbrace \pm p^i s_1 p^j |
-1 \leq i,j\leq 1\rbrace $. The matrix $p$ is invertible
and $s_1\not=0$ so
$p^i s_1 p^j\not=0$ for $-1 \leq i,j \leq 1$.
Therefore $m\not=0$ and the claim is proved.
Let $\kappa $ denote the $(1,1)$-entry of the matrix $\eta $.
Then $\eta=\kappa I$ since $\eta \in Z(\Lambda )$.
In the equation
$\eta_1 \eta_2\cdots \eta_n=\kappa I$ take the determinant of
each side and use
Lemma~\ref{lem:4}$(i)$
to get $(-\lambda)^\ell = \kappa^2$.
Therefore $\ell $ is even and
$\kappa=\pm \lambda^{\ell/2}$.
Now $\eta =\pm \lambda^{\ell/2} I$, so
the above polynomial $f$ has degree $\ell/2$.
But $\ell \not=0$ so $\ell > \ell/2$
for a contradiction.
the result follows.
\end{proof}

\begin{thebibliography}{99}

\footnotesize\itemsep=0pt
\bibitem{alpern}
Alperin R.C.,
Notes: ${\rm  {PSL}}_2(\mathbb Z)= \mathbb Z_2 \star \mathbb Z_3$,
\href{http://dx.doi.org/10.2307/2324963}{{\it Amer. Math. Monthly}} {\bf 100} (1993), 385--386.
\end{thebibliography}
\end{document}
